\section{Open Questions 与长期方向}
\subsection{可解释性与监管}
Weston 认为 ``the transformer itself does reasoning'' 但又提醒我们 ``interpreting the neural network isn't easy''（SRT 3204）。因此，长期方向之一是构建可追溯的 reasoning pipeline，让 verified response 除了 final answer 外还能输出 trace 以供监管使用。

\begin{table}[H]
\centering
\begin{tabular}{p{5cm}p{10cm}}
\toprule
问题 & 描述 \\
\midrule
解释性缺口 & 如何从 verified response 中提取 reasoning trace 以供人类审查？ \\
责任归属 & meta judge 判错时，如何追溯是 reward model、judge 还是 prompt 操作失误？ \\
多代理信号 & 在多 agent 场景（bot ↔ judge ↔ user）中，如何保证 reward signal 不被 adversarial exploitation？ \\
\bottomrule
\end{tabular}
\caption{Open Questions 供后续 alignment 与审计使用。}
\end{table}

\subsection{与外部系统的交互}
未来的自我提升模型需要 ``Reason by actually interacting with people in the world, internet or itself''（SRT 3282-3284）。这包括：
\begin{itemize}
    \item 与其他 agent 联动（share verified responses, cross judge）；
    \item 通过 API / internet 搜集 additional context，再在 Self-Instruct 里自我验证；
    \item 在发现 hallucination 时自动寻求 external agent 干预（self-aware behaviour）。
\end{itemize}

\subsection{本章小结}
Open Questions 提供了 alignment 与监管的方向；与外部世界的交互则是搭建更自主、更自知、也更安全的 self-improving 体系的下一步。

